\documentclass[12pt]{article}

\setlength{\parindent}{0pt}
\setcounter{secnumdepth}{3}
\usepackage[margin=1in]{geometry}

% packages
\usepackage{amsmath, amsfonts, amsthm, amssymb}
\usepackage{graphicx, float}
\usepackage{mathtools}
\usepackage{titlesec}
\usepackage{interval}
\usepackage{hyperref}
\usepackage{siunitx}
\usepackage{titling}
\usepackage{vwcol}
\usepackage{setspace}
\usepackage{empheq}
\usepackage{cancel}
\usepackage{esdiff}
\usepackage{multicol}
\usepackage{mdframed}
\usepackage{esdiff}
\usepackage{tikzsymbols}
\usepackage{multicol}
\usepackage{tikz}
\usepackage{varwidth}
\usepackage{pgfplots}
\usepackage{fancyhdr}
\usepackage{enumitem}
\usepackage{tcolorbox}

\intervalconfig {
	soft open fences
}

% header
\pagestyle{fancy}
\fancyhf{}
\lhead{Ethan Fan}
\rhead{PCH}
\cfoot{\thepage}

% section format
\titleformat{\section}
  {\large \bfseries}
  {}
  {0em}
  {}
\titlespacing*{\section}{0pt}{0.5em}{0.3em}

\titleformat{\subsection}[runin]
  {\bfseries}
  {}
  {0em}
  {}
\titlespacing*{\subsection}{0pt}{0pt}{0em}

\titleformat{\subsubsection}[runin]
  {\itshape}
  {(\roman{subsubsection})}
  {0em}
  {}

% commands
\newcommand{\alignedintertext}[1]{%
  \noalign{%
    \vskip\belowdisplayshortskip
    \vtop{\hsize=\linewidth#1\par
    \expandafter}%
    \expandafter\prevdepth\the\prevdepth
  }%
}

\newcommand*{\paren}[1]{\ensuremath\left(#1\right)}
\newcommand*{\limit}[2][x]{\ensuremath{\displaystyle\lim_{#1\to#2}}}
\newcommand*{\Limit}[3][x]{\ensuremath{\displaystyle\lim_{#1\to#2}\left[#3\right]}}
\newcommand*{\deriv}[1][x]{\ensuremath{\dfrac{\mathrm{d}}{\mathrm{d}#1}}}
\newcommand*{\Deriv}[2][x]{\ensuremath{\dfrac{\mathrm{d}}{\mathrm{d}#1}\left[#2\right]}}
\newcommand{\dydx}{\dfrac{dy}{dx}}
\newcommand{\dydt}{\dfrac{dy}{dt}}
\newcommand{\dxdt}{\dfrac{dx}{dt}}
\newcommand*{\iinteg}[2][x]{\ensuremath{\displaystyle\int #2\;\mathrm{d}#1}}
\newcommand*{\dinteg}[4][x]{\ensuremath{\displaystyle\int_{#2}^{#3}#4\;\mathrm{d}#1}}
\newcommand*{\abs}[1]{\ensuremath{\left|#1\right|}}
\newcommand*{\eps}{\varepsilon}
\newcommand*{\real}{\ensuremath\mathbb{R}}
\newcommand*{\integer}{\ensuremath\mathbb{Z}}
\newcommand*{\posint}{\ensuremath\mathbb{Z^+}}
\newcommand*{\cbrt}[1]{\ensuremath{\sqrt[3]{#1}}}
\newcommand*{\lheqzero}{\ensuremath{\underset{\text{L'H}}{\overset{\left[\frac00\right]}{=}}}}
\newcommand*{\lheqinfty}{\ensuremath{\underset{\text{L'H}}{\overset{\left[\frac{\infty}{\infty}\right]}{=}}}}
\newcommand{\tor}{\text{ or }}
\newcommand{\floor}[1]{\lfloor #1 \rfloor}

\newcommand{\vem}{\vspace{0.5em}}
\newcommand{\example}[1]{\section{Example #1}}
\newcommand{\problem}[1]{\section{Problem #1}}
\newcommand{\subproblem}[1]{\subsection{(#1)} \,}

\DeclareMathOperator{\dne}{DNE}
\DeclareMathOperator{\sgn}{sgn}

\newcommand{\definition}[1]{\begin{tcolorbox}[colback=red!5!white,colframe=red!75!black,parbox=false] #1 \end{tcolorbox}}

\newcommand{\theorem}[2]{\begin{tcolorbox}[title={#1},colback=blue!5!white,colframe=blue!75!black,parbox=false] #2 \end{tcolorbox}}

\allowdisplaybreaks
% \postdisplaypenalty=100000

% document
\begin{document}

\vspace*{0cm}

% opening
\begin{center}
    {\Large \bfseries Problem Set 77} \\[0.5cm]
    {\large Collection of Problems} \\[0.35cm]
    {\normalsize May 25, 2026}
\end{center}

% problems

\problem{3}
Let $L=\lim\limits_{x\to 0} \dfrac{a-\sqrt{a-x^2}-\dfrac{x^2}{4}}{x^4}$ be a finite limit. Find the values of $a(a>0)$ and $L$.
\begin{gather*}
    L=\lim\limits_{x\to 0} \dfrac{a-\sqrt{a-x^2}-\dfrac{x^2}{4}}{x^4} \quad \left[\frac{0}{0}\right]\\
    \begin{cases}
        \lim\limits_{x\to 0} a-\sqrt{a-x^2}-\dfrac{x^2}{4}=0\\
        \lim\limits_{x\to 0} x^4 =0
    \end{cases}\\
    L \lheqzero \lim\limits_{x\to 0}\frac{\frac{x}{\sqrt{a-x^2}}-\frac{x}{2}}{4x^3} = \lim\limits_{x\to 0}\frac{\frac{1}{\sqrt{a-x^2}}-\frac{1}{2}}{4x^2}\\
    \begin{cases}
        \lim\limits_{x\to 0} 4x^2=0\\
        \lim\limits_{x\to 0} \dfrac{1}{\sqrt{a-x^2}}-\dfrac{1}{2}=0 \implies a=2
    \end{cases}\\
    L\lheqzero \lim\limits_{x\to 0}\frac{-\frac{1}{2}\cdot \frac{-2x}{(2-x^2)^{3/2}}}{8x}= \lim\limits_{x\to 0} \frac{1}{8\cdot 2^{3/2}}=\frac{\sqrt{2}}{32}\\
    \boxed{a=2, L=\frac{\sqrt{2}}{32}}
\end{gather*}

\problem{4}
Let $f(x)=(1-x)^2(\sin x)^2+x^2$ for all $x\in \mathbb{R}$. Consider the statements: \vem

P: There exists some $x\in \mathbb{R}$ such that
\[
f(x)+2x=2(1+x^2)
\]
Q: There exists some $x\in \mathbb{R}$ such that
\[
2f(x)+1=2x(1+x)
\]
Then \vem

(A) both P and Q are true \quad (B) P is true and Q is false \vem

(C) P is false and Q is true \quad (D) both P and Q are false\vem

Justify your answer.
\begin{gather*}
    $P$: f(x)=2x^2-2x+2\\
    $Q$:f(x)=x^2+x-\frac{1}{2}
\end{gather*}
For P:
\begin{align*}
    (1-x)^2(\sin x)^2+x^2&=2x^2-2x+2\\
    (1-x)^2\sin ^2x&=x^2-2x+2\\
    (1-x)^2\sin ^2x&=(1-x)^2+1\\
    (1-x)^2\cos^2x&=-1
\end{align*}
As no real $x$ satisfies the above equality, P is false. For Q:
\begin{align*}
    h(x)&=2f(x)+1-2x(1+x)\\
    h(0)&=2f(0)+1-0=1\\
    h(1)&=2f(1)+1-2\cdot 2=-2
\end{align*}
By the Intermediate Value Theorem, $h(x)$ holds a solution on $(0,1)$. Hence, $Q$ is true. 
\[
\boxed{$(C) P is false and Q is true$}
\]

\problem{5}
Discuss continuity and differentiability
\[
f(x)=\begin{cases}
    -x-\dfrac{\pi}{2} &$if $ x\le -\dfrac{\pi}{2}\\
    -\cos x &$if $-\dfrac{\pi}{2}<x\le 0\\
    x-1 &$if $ 0<x\le 1\\
    \ln x &$if $ x>1
\end{cases}
\]
\[
f'(x)=\begin{cases}
    -1 &$if $ x\le -\dfrac{\pi}{2}\\
    \sin x &$if $-\dfrac{\pi}{2}<x\le 0\\
    1 &$if $ 0<x\le 1\\
    \dfrac{1}{x} &$if $ x>1
\end{cases}
\]
At $x=-\dfrac{\pi}{2}$:
\begin{gather*}
    f(-\frac{\pi}{2}^-)=0,f(-\frac{\pi}{2}^+)=0,f(-\frac{\pi}{2})=0\\
    f(-\frac{\pi}{2}^-)=f(-\frac{\pi}{2}^+)=f(-\frac{\pi}{2})\\
    f'(-\frac{\pi}{2}^-)=-1,  f'(-\frac{\pi}{2}^+)=-1\\
    f'(-\frac{\pi}{2}^-)= f'(-\frac{\pi}{2}^+) \implies f'(-\frac{\pi}{2})=-1
\end{gather*}
$f(x)$ is continuous and differentiable at $x=-\dfrac{\pi}{2}$. At $x=0$:
\begin{gather*}
    f(0^-)=-1,f(0^+)=-1,f(0)=-1\\
    f(0^-)=f(0^+)=f(0)\\
    f'(0^-)=0, f'(0^+)=1\\
    f'(0^-)\ne f'(0^+) \implies f'(0)\dne 
\end{gather*}
$f(x)$ is continuous but not differentiable at $x=0$. At $x=1$:
\begin{gather*}
    f(1^-)=0, f(1^+)=0, f(1)=0\\
    f(1^-)=f(1^+)=f(1)\\
    f'(1^-)=1, f'(1^+)=1\\
    f'(1^-)=f'(1^+) \implies f'(1)=1
\end{gather*}
$f(x)$ is continuous and differentiable at $x=1$.

\problem{6}
Find the domain of the derivative function
\[
f(x)=\begin{cases}
    \arctan x & $if $ |x|\le 1\\
    \dfrac{1}{2}(|x|-1) &$if $ |x|>1
\end{cases}
\]
\[
    f'(x)=\begin{cases}
        \dfrac{1}{1+x^2} & $if $ |x|\le 1\\
        \dfrac{1}{2}\sgn (x) &$if $ |x|>1
    \end{cases}
\]
At $x=-1$:
\begin{gather*}
    f'(-1^-)=-\frac{1}{2},f'(-1^+)=\frac{1}{2}\\
    f'(-1^-)=f'(-1^+)\implies f'(-1)\dne
\end{gather*}
$f(x)$ is not differentiable at $x=-1$. At $x=1$:
\begin{gather*}
    f'(1^-)=\frac{1}{2}, f'(1^+)=\frac{1}{2}\\
    f'(1^-)=f'(1^+)\implies f'(1)=\frac{1}{2}
\end{gather*}
$f(x)$ is differentiable at $x=1$.
\[
\boxed{D_{f'} =(-\infty ,-1) \cup (-1,\infty)}
\]

\problem{7}
Let $g(x)=\ln f(x)$ where $f(x)$ is twice differentiable positive function on $(0,\infty)$ such that $f(x+1)=xf(x)$. Find
\[
g''\paren{x+\frac{1}{2}}-g''\paren{\frac{1}{2}}
\]
(in terms of $N$), where $x=1,2,3,\dots ,N.$
\begin{gather*}
    \ln f(x+1)=\ln x+\ln f(x)\\
    g(x+1)-g(x)=\ln x\\
    g'(x+1)-g'(x)=\frac{1}{x}\\
    g''(x+1)-g''(x)=-\frac{1}{x^2}\\ 
    x=\frac{1}{2} \implies g''(\frac{3}{2})-g''(\frac{1}{2})=-4\\
    x=\frac{3}{2} \implies g''(\frac{5}{2})-g''(\frac{3}{2})= -\frac{4}{9}\\
    x=\frac{5}{2} \implies g''(\frac{7}{2})-g''(\frac{5}{2})=-\frac{4}{25}\\
    \boxed{g''(N+\frac{1}{2})-g''(\frac{1}{2})=-4 \sum_{k=1}^N\frac{1}{(2k-1)^2}}
\end{gather*}

\problem{8}
Let $f(\theta )=\sin \paren{\arctan \paren{\dfrac{\sin \theta}{\sqrt{\cos 3\theta}}}}$, where $-\dfrac{\pi}{4}<\theta < \dfrac{\pi}{4}$. Find $\dfrac{d}{d(\tan \theta)}(f(\theta ))$.
\begin{align*}
    &\dfrac{d}{d(\tan \theta)}(f(\theta ))\\
    &= \frac{\frac{df}{d\theta}}{\frac{d(\tan \theta)}{d\theta}}\\
    &= \frac{\cos \paren{\arctan \paren{\frac{\sin \theta}{\sqrt{\cos 3\theta}}}} \paren{\frac{1}{1+ \frac{\sin ^2\theta}{\cos 3\theta}}} (\frac{\sin \theta}{\sqrt{\cos 3\theta}})'}{\sec ^2\theta}\\
    &=\cos \paren{\arctan \paren{\frac{\sin \theta}{\sqrt{\cos 3\theta}}}} \paren{\frac{1}{\sec^2 \theta(1+ \frac{\sin ^2\theta}{\cos 3\theta})}} \frac{\cos \theta \sqrt{\cos 3\theta}+\sin \theta \frac{3\sin 3\theta}{2\sqrt{\cos 3\theta}}}{\cos 3\theta}\\
    &=\frac{\sqrt{\cos 3\theta}}{\sqrt{\cos 3\theta + \sin^2 \theta}}\paren{\frac{\cos^2 \theta \cos 3\theta}{\cos 3\theta+ \sin ^2\theta}} \frac{\cos \theta \sqrt{\cos 3\theta}+\sin \theta \frac{3\sin 3\theta}{2\sqrt{\cos 3\theta}}}{\cos 3\theta}\\
    &= \boxed{\frac{\cos^2 \theta (2\cos\theta \cos 3 \theta + 3\sin \theta \sin 3 \theta)}{2(\cos 3\theta + \sin^2 \theta)^\frac{3}{2}}}
\end{align*}


\end{document}

